% !TEX program = xelatex
% =============================================================================
%  中文默认数学建模论文 LaTeX 模板
%  对应 Typst 模板 zh/default/main.typ 的全部功能。
%
%  字体依赖（正文宋体 / 标题黑体 / 强调楷体）：
%    - 规范字体：SimSun（宋体）、SimHei（黑体）、KaiTi（楷体）。
%    - ctex 自动选择当前系统字体集；跨平台使用前编译确认字体可用。
%    - 西文正文：Times New Roman；代码等宽：Menlo（缺失时回退到默认字体）。
%  编译方式：xelatex main.tex （需运行两遍以生成目录）
% =============================================================================
\documentclass[12pt, a4paper]{ctexart}

% =====================================================================
%  本族只写「赛事/族差异」。通用排版规范全在 _base/ 里（复制模板时必须把
%    templates/_base/ 一并复制到 paper/_base/，否则编译不过）。
%      _base/preamble.tex   页面/字体/行距/标题骨架/caption/表格/浮动体/代码/断行
%      _base/macros.tex     \papertitle \abstractcn \referencescn \appendixAcn
%                           \threelinetable \paperfigure \paperfigurepair…
%      _base/math-style.tex 数学字体与环境
%    通用排版规范（页边距、\linespread、graphicx、caption 统一、刚性行距）
%    全部由 _base/ 提供，本文件不再自带导言区。
% =====================================================================
\input{_base/preamble.tex}
\input{_base/macros.tex}

% ---------- 结果登记值（由 `python -m lib.result_contract export` 生成，禁止手工编辑） ----------
% docs/RESULT_CONTRACT.md 是全链通用要求：正文数值一律写 \ResultValue{指标ID}，
% 不写裸数字。\ResultValue 未定义（result-values.tex 尚未生成）时直接排会触发
% `! Undefined control sequence`。
% \detokenize 的理由：指标 ID 带下划线，直接排文本模式会触发 Missing $。
\IfFileExists{result-values.tex}{\input{result-values.tex}}{}
\providecommand{\ResultValue}[1]{\textbf{[未生成 result-values.tex：\detokenize{#1}]}}

% ---------- 族差异 ①：编号用阿拉伯数字（_base 默认是国赛的中文数字「一、」） ----------
\ctexset{
  section/number       = \arabic{section},
  subsection/number    = \arabic{section}.\arabic{subsection},
  subsubsection/number = \arabic{section}.\arabic{subsection}.\arabic{subsubsection},
}
% ---------- 族差异 ②：标题字号按本族规范 16/14/12pt（_base 默认 17/15/13） ----------
\titleformat{\section}
  {\centering\fontsize{16pt}{19.2pt}\bfseries}
  {\arabic{section}}{1em}{}
\titleformat{\subsection}
  {\fontsize{14pt}{16.8pt}\bfseries}
  {\arabic{section}.\arabic{subsection}}{1em}{}
\titleformat{\subsubsection}
  {\fontsize{12pt}{14.4pt}\bfseries}
  {\arabic{section}.\arabic{subsection}.\arabic{subsubsection}}{1em}{}
\titlespacing{\section}       {0pt}{1.8em}{1.0em}
\titlespacing{\subsection}    {0pt}{1.25em}{0.65em}
\titlespacing{\subsubsection} {0pt}{0.9em}{0.4em}

% ---------- 目录（tocloft 点线引导） ----------
\usepackage{tocloft}
\renewcommand{\cftsecfont}{\heiti\bfseries}
\renewcommand{\cftsecpagefont}{\heiti\bfseries}
\renewcommand{\cftsecleader}{\cftdotfill{\cftdotsep}}
\renewcommand{\cftsubsecleader}{\cftdotfill{\cftdotsep}}
\renewcommand{\cftsubsubsecleader}{\cftdotfill{\cftdotsep}}

% ---------- 页眉页脚（fancyhdr） ----------
% 对应 Typst：页眉“数学建模竞赛论文”居中 + 下划线，页脚页码居中
\usepackage{fancyhdr}
\pagestyle{fancy}
\fancyhf{}
\fancyhead[C]{\fontsize{10.5pt}{12.6pt}数学建模竞赛论文}
\fancyfoot[C]{\thepage}
\renewcommand{\headrulewidth}{0.4pt}

% ---------- 族差异 ④：目录页 ----------
% \papertitle / \abstractcn / \referencescn / \appendixAcn / \threelinetable /
% \paperfigure 等全部在 _base/macros.tex 里，本族不再重复定义 —— 重复定义会直接
% 撞名报错（`Command \papertitle already defined`）。
\makeatletter
\newcommand{\tocpage}{%
  \setcounter{page}{1}%
  \begin{center}
    {\fontsize{16pt}{19.2pt}\bfseries 目录}
  \end{center}
  \vspace{1em}
  \@starttoc{toc}%
  \newpage
}
\makeatother

\begin{document}

% 标题页（不编号、不显示页码）
\vspace*{1.85cm}
\papertitle{论文标题}

\vspace{2.25cm}

% 摘要页（不编号）
\abstractcn{%
  摘要内容：问题概述 + 每个子问题的方法和数值结果 + 结论%
}{%
  关键词1；关键词2；关键词3%
}

% 目录页（页码从 1 开始，带页眉）
\tocpage

\section{问题重述}
% 依据 docs/WRITING_QUALITY.md，从本题提炼背景与每问任务；按实际问数组织。
% 每问说明决策对象、变化条件和交付物。详细参数放数据章节。
% 区分题面要求的附件与实际收到的文件；缺失材料明确披露。

\section{模型假设}

为简化问题，本文做出以下基本假设：

\begin{enumerate}
  \item 假设题目所给数据真实可靠，不存在系统性的测量误差；
  \item 假设各因素之间的交互作用可以忽略不计；
  \item 假设系统在短期内处于稳定状态；
  \item 假设所有变量连续可微；
  \item 假设缺失值为随机缺失，可采用插补方法处理。
\end{enumerate}

\section{符号说明}
% 待填：本节内容由 9Paper-writing 按 ANALYSIS_MODELING_REPORT.md / RESULTS_REPORT.md 撰写。
%
% 形制与国赛模板 `zh/cumcm-latex/sections/4_symbols.tex` 一致，照那份写：
% 三列「符号 / 含义 / 单位」，单位单独成列（不写进含义的括号里）、左/中/右对齐、
% 用 longtable 而非浮动体、
% 表标题放在表外。完整骨架与逐条注释见那份文件。

\section{问题一的模型建立与求解}
\subsection{数据预处理}
\subsection{模型建立}
\subsection{求解结果}
% 待填：本节内容由 9Paper-writing 按 ANALYSIS_MODELING_REPORT.md / RESULTS_REPORT.md 撰写。

\section{问题二的模型建立与求解}
\subsection{特征选择}
\subsection{模型优化}
% 待填：本节内容由 9Paper-writing 按 ANALYSIS_MODELING_REPORT.md / RESULTS_REPORT.md 撰写。

\section{问题三的模型建立与求解}
\subsection{优化模型}
\subsection{算法求解}
\subsection{结果}
% 待填：本节内容由 9Paper-writing 按 ANALYSIS_MODELING_REPORT.md / RESULTS_REPORT.md 撰写。

\section{敏感性分析}
\subsection{参数敏感性}
% 待填：本节内容由 9Paper-writing 按 ANALYSIS_MODELING_REPORT.md / RESULTS_REPORT.md 撰写。

\section{模型评价与推广}

\subsection{模型优点}

\begin{enumerate}
  \item 理论基础扎实，推导严谨；
  \item 多模型对比确保最优性；
  \item 可解释性强，计算效率高。
\end{enumerate}

\subsection{模型缺点}

\begin{enumerate}
  \item 对数据质量有一定要求；
  \item 未考虑时间动态性。
\end{enumerate}

\subsection{推广}

可推广至金融风控、工业生产、智能交通和能源管理等领域。


\newpage
\referencescn
\newpage
% 附录（正文的最后一页）：只放「支撑材料文件清单」表，由 14Layout-and-format 据
% reports/SUBMISSION_MANIFEST.json 生成，必须与打包出来的产物逐项一致。
% 详细的补充推导/图表单独编译成 `附录A.pdf`（工程在 `paper_appendix/`，模板见 `appendix/`）。
% 这里没有「附录 B 完整源程序」—— 源码不再排进论文，改以散装 .py 提交。
\appendixAcn[支撑材料清单]{A1_materials}

\end{document}
